% Reconstructed/authored after executor recovery, 2026-11-20.
% Uses the physically recovered complete A.115 and the audited A.156/A.159 arguments.
\Needspace{10\baselineskip}
\subsection{A.160. The PL--definable comparison over every real closed field}

\begin{saresult}{A.160}{Proved under stated hypotheses}{All-field weak equivalence}
Let $M$ be any fixed o-minimal expansion of a real closed field,
and let $B$ be a closed compact PL manifold over $M$. Then
\[
 P_{\mathrm{PL}}(B)\longrightarrow P_{\mathrm{def}}(B)
 \tag{160.1}
\]
is a weak equivalence. The statement compares finite jointly
PL parameter families with finite jointly definable parameter
families. It does not identify either space with the whole
ordinary TOP concordance space in exceptional dimensions,
nor does it by itself compute their homotopy type.
\end{saresult}

\begin{proof}
Higher surjectivity is A.159 and the component bijection is
A.113. We prove injectivity, retaining the exact parameter
projection throughout. Replace $B$ by its standard abstract
finite PL model as in A.159 and write $F=\mathbb R_{\mathrm{alg}}$
for the ordered real closed subfield of real algebraic numbers
in $M$. The fixed model and its two end markings are defined
over $F$.

\textbf{The recovered parameter-preserving field comparison.}\quad
We use the full statement and proof of A.115. This is not an
appeal to field invariance of ordinary, unparameterized PL
isotopies. Its proof simultaneously subdivides the one-way
chain
\[
 N_M\times P_M\xrightarrow{h}N_M\times P_M
                  \xrightarrow{\operatorname{pr}_P}P_M,
\]
retains fixed product-cell carriers, and transfers the finite
vertex configuration to $F$. The matching-vertex maps $u,v$
cover the SAME base map $\beta$. Thus
$(\mathrm{id}_N\times\beta^{-1})u$ and
$(\mathrm{id}_N\times\beta^{-1})v$ preserve the parameter
pointwise and preserve every fixed product cell.
The finite total-product cone construction contracts them
with their agreement on the bottom and parameter boundary
retained. The same construction removes the base conjugation
by $\beta$. This produces an actual jointly PL concordance
homotopy to an $F$-coefficient representative. No polygonal
approximation to a coefficient-space path is used.

Let $h$ be a based jointly PL $k$-cube over $M$, $k\ge1$,
which is nullhomotopic through based definable concordances.
By A.115, choose a based jointly PL homotopy from $h$ to
an $F$-PL cube $h_0$. Concatenate the reversed homotopy
with the given definable nullhomotopy of $h$. We obtain a
definable filling $U$ of the full boundary data on a
$(k+1)$-cube $P$: $h_0$ on the initial face and the identity
on the other faces. These full boundary data are jointly PL
over $F$ and agree exactly on their intersections.

We now prove that such a filling with $F$-PL boundary admits
a PL filling over $M$ with exactly the same boundary.

\textbf{Standard bounds for the boundary factors.}\quad
Use the fixed $F$-PL cubical model $Q\subset C=[0,1]^n$ of
A.154, and write $U^Q$ for the conjugated family.
Its boundary restriction has a PL vector extension
\[
 E_\partial:C\times\partial P\longrightarrow C
\]
over $F$: take a finite compatible triangulation over $F$,
leave the prescribed vertex values unchanged, assign a
fixed rational point of $C$ to the other vertices, and
extend affinely. This is a vector extension and assumes
no extension by homeomorphisms. Paste it with $U^Q$ on
$Q\times P$ and use the full-subcomplex extension of A.154
to obtain a continuous definable vector map $E$ on $C\times P$.

Run the PL-normalized coordinate clocks of A.156, starting
with $(\mathrm{id},E)$. At each stage all component values
lie in a fixed unit cube. If their number is $d$ and the
uniform monotonicity count is $N$, the raw clock height is
at most $1+dN$. Choose its trapezoid normalization parameter
$D$ to be a STANDARD INTEGER larger than $\max(3,1+dN)$.
The counts $d,N$ are ordinary finite integers even though
they may depend on the input over $M$.

On the parameter boundary the initial data are PL over $F$.
The finite maximum formula for line variation is semilinear
over $F$ when the input is PL over $F$; continuity makes
the resulting variation function PL over $F$. The
normalization $T_D$ of A.156 has coefficients in $F$,
and the inverse of the resulting parameter-preserving PL
clock has coefficients in $F$ as well. Induction shows
that every boundary clock, and the two final boundary
factors, remain jointly PL over $F$.
Hence the factorization
\[
 U_p=s_pr_p^{-1}
 \tag{160.2}
\]
has uniformly spatially Lipschitz factors with standard
bounds, and its boundary factors have STANDARD forward
parameter Lipschitz bounds. The latter assertion follows
from their finite PL descriptors over the Archimedean field
$F$; it is not a rounding of an arbitrary element of $M$.

The factors individually fix the bottom. Indeed the bottom
is a union of coordinate faces of $Q$, and $E=\mathrm{id}$
there. On each free coordinate line in such a face the
evolving vector is $(\mathrm{id},\mathrm{id})$, so its raw
clock is $3t$. The normalization of A.156 is exactly the
identity at height $3$. Coordinates fixed by that face
remain fixed endpoints. Induction gives $r=s=\mathrm{id}$
on the bottom. Both factors preserve the top setwise.

\textbf{Exact relative graph regularization.}\quad
Form the two tagged graph family as in A.159. It has a
standard spatial bound and its parameter-boundary graph
family has a standard Hausdorff bound, by the preceding
$F$-PL assertion. A.158 gives a standard compact Lipschitz
resolution. Add the exact boundary copy of A.157.
Its tag coordinates, its projection, and its nonconstant
pulled boundary family all have standard bounds.
The finite join estimate and bounded vector clocks
therefore give a map $g:P\to P$ fixing $\partial P$
pointwise and standard Hausdorff control for both pulled
graph families. As in (159.10), the common standard
spatial bound converts this to a standard forward
parameter bound for
\[
 \bar r=r\circ g,\qquad \bar s=s\circ g.
\]
Their boundary restrictions are still exactly $r,s$,
and their bottom restrictions are still the identity.

\textbf{The exact PL endpoint and its boundary collar.}\quad
Apply the extension-only relative argument of A.156 to
each factor. Its only Archimedean use was the choice of
a finite Hudson grid against the forward parameter bound.
That bound is now standard, so A.159 supplies the fixed
finite target data over $F$. The contraction is solved
over $M$ by the residual-minimum argument.

The PL roof normalization of A.156 gives a jointly PL
homotopy track on the parameter boundary. Its inverse
formula evaluates only the given PL boundary inverse.
The target prisms and all boundary source prisms are PL;
the common marked ambient triangulation preserves these
boundary cells setwise. Thus the ORIGINAL boundary map
has the required correspondence after conjugation.
Prescribe it, together with the bottom identity, first.
The total-product cone induction extends it over the
remaining finite prisms. Lower parameter faces may carry
spatial subdivisions; their prisms lie in the already
assigned lower-dimensional product boundary. No PL filling
is assumed at this step.

A fixed PL collar of an inner parameter cube, carrying
the reversed PL boundary track, restores the initial
boundary exactly. All these operations are finite PL
operations over $M$ and require no further metric mesh.
Consequently we obtain jointly PL bottom-fixed families
$R,S$ with the boundary values of $r,s$, respectively.
Their quotient
\[
 W_p=S_pR_p^{-1}
 \tag{160.3}
\]
is a PL filling of $U|_{\partial P}$ with exactly its
original end markings and parameter-boundary values.

In particular $h_0$ is PL nullhomotopic over $M$.
Concatenate this nullhomotopy with the PL homotopy
$h\simeq h_0$ supplied by A.115. The result is a PL
nullhomotopy of $h$, proving injectivity on $\pi_k$.
All concatenations use fixed affine subdivisions of the
time interval, so are jointly PL where asserted.
Together with A.159 and A.113 this proves (160.1).
Translation by PL representatives handles all base
components. Empty or disconnected bases are treated
componentwise as in the preceding results.
\end{proof}

